% Generated draft: requires mathematical and source review.
\documentclass[11pt]{article}
\usepackage[margin=25mm]{geometry}
\usepackage{fontspec}
\setmainfont{Latin Modern Roman}
\usepackage{amsmath,amssymb,mathtools,unicode-math}
\setmathfont{Latin Modern Math}
\usepackage{enumitem,needspace,xurl,hyperref}
\setlength{\parindent}{0pt}
\setlength{\parskip}{5pt}
\setlength{\emergencystretch}{3em}
\newcommand{\N}{\mathbb{N}}
\newcommand{\Z}{\mathbb{Z}}
\newcommand{\Q}{\mathbb{Q}}
\newcommand{\R}{\mathbb{R}}
\newcommand{\C}{\mathbb{C}}
\newcommand{\F}{\mathbb{F}}
\newcommand{\A}{\mathbb{A}}
\newcommand{\PP}{\mathbb{P}}
\newcommand{\E}{\mathbb{E}}
\newcommand{\eps}{\varepsilon}
\newcommand{\dd}{\,\mathrm{d}}
\begin{document}
\section*{UnsolvedMath ID: 30000190}
\textbf{Source problem code:} catalogue entry
\subsection{Definitions and mathematical statement}
Let $\mathcal A_g^F$ denote the first Voronoi, or perfect-cone, toroidal compactification of the moduli space of principally polarized abelian varieties of dimension $g$, as identified by companion permanent ID $30000189$. Let $E$ be the extended rank-$g$ Hodge bundle. In the standard notation, $M$ is the Hodge determinant divisor class $c_1(\det E)$ and $D$ is the boundary divisor. These conventions must be matched to the report's stack/coarse-space and base-field setting.

The source contains two questions:
\begin{enumerate}
\item Is the divisor $12M-D$ semi-ample on $\mathcal A_g^F$? For a rational divisor, this means that some positive integral Cartier multiple $m(12M-D)$ has globally generated associated line bundle.
\item Ekedahl's question: which ample line bundles $L$ on $\mathcal A_g^F$ have $E\otimes L$ ample? Vector-bundle ampleness means that the tautological quotient line bundle $\mathcal O_{\mathbb P(E\otimes L)}(1)$ is ample.
\end{enumerate}
The second question asks which twists themselves are ample, not merely eventual ampleness after taking an unspecified large power. The two questions are separate: semiampleness of one determinant/boundary combination does not supply a classification of ample twists of a rank-$g$ bundle.

The extract does not specify characteristic, genus range, level structure, or whether bundles are interpreted on the moduli stack, a cover, or the coarse compactification. These choices matter because the Hodge bundle need not be an honest vector bundle on the unqualified coarse space. The displayed notation uses the usual determinant and boundary interpretation without asserting a descent convention absent from the source. In particular, the $6\le g\le10$ range of the neighboring flip problem is not imported into these two questions.

The literature assessment checked on 21 August 2026 reports that $12M-D$ is nef, is semi-ample in positive characteristic, and has an identified exceptional locus in all characteristics. It distinguishes those facts from characteristic-zero semiampleness and from the separate Ekedahl ampleness question, which it does not present as jointly settled. Knowledge of an exceptional locus or nefness alone is not a semiampleness proof. These are supplied dated partial-progress claims requiring independent verification.
\subsection{Short English statement}
On the perfect-cone abelian compactification, is $12M-D$ semi-ample, and which ample line-bundle twists make the Hodge bundle ample? The supplied progress is characteristic-sensitive and does not combine the two questions into one solved assertion.
\subsection{Sources}
\begin{itemize}
\item \emph{Komplexe Algebraische Geometrie} (2005), Oberwolfach Reports, DOI 10.4171/owr/2005/07, pp. 407--463, questions (5) and (6), the latter attributed to Ekedahl, \url{https://doi.org/10.4171/owr/2005/07}.
\item Literature assessment checked 21 August 2026: Shepherd-Barron, ``exceptional locus in the perfect compactification,'' \url{https://arxiv.org/abs/1604.05954}; original report, \url{https://doi.org/10.4171/owr/2005/07}.
\item Companion local record, permanent ID $30000189$, source code \texttt{OWR-788-007}, identifies $\mathcal A_g^F$ as the first Voronoi compactification in the neighboring explicit-flip question.
\item Local UnsolvedMath registry, permanent ID $30000190$, exact source code \texttt{catalogue entry}, created 25 August 2026 and updated 16 September 2026. The clean statement is labelled corrected and verified, with compactification notation and the Hodge-bundle identity recovered.
\end{itemize}
\end{document}
